\section{Method}
\label{sec:method}

\ingram{} is a layer inserted between two blocks of a sequence model, and each insertion is a
\emph{placement}. At every position the layer writes to and reads from a memory whose entries are
addressed by a hash of the preceding $N$ token ids, the \emph{cue}; positions interact only
through this memory. Figure~\ref{fig:overview}b shows one write and one read for one cue length.

\subsection{Cues, slots and table rows}
\label{sec:method-cues}

Let $w_1, \dots, w_L$ be the token ids of a context, and let $x_t \in \mathbb{R}^d$ be the hidden
state at position $t$ where the layer is inserted. For a cue length $N$, the cue at position
$t \ge N$ is the $N$-gram $c = (w_{t-N+1}, \dots, w_t)$ that ends there. A polynomial hash maps it
to a 61-bit address,
\begin{equation}
a_N(c) = \Bigl( \sum_{i=0}^{N-1} (w_{t-i} + 1)\, p_N^{\,i} \bmod (2^{61} - 1) \Bigr) \oplus s_N ,
\label{eq:hash}
\end{equation}
with a prime $p_N$ and a constant $s_N$ for each cue length. We use six cue lengths, $N \in \{1,
2, 4, 8, 16, 32\}$, and each has its own memory.

A cue addresses two vectors of width $d_m = 64$, a slot and a table row. The \emph{slot} $S_c$
belongs to the current context: it is created at zero when the context first writes to $c$ and is
keyed by the full address, and two different cues therefore share a slot only if their 61-bit
addresses collide. The \emph{table row} $R_c$ is a parameter that all contexts share, one row of a
learned table with $2^{16}$ rows per cue length, selected by a second hash of the address (the
splitmix64 finalizer) reduced modulo $2^{16}$. Different cues can
share a row, as in hashed $N$-gram embeddings \citep{engram2026}.

\subsection{Write and read}
\label{sec:method-rw}

At every position $p$ and for each cue length, the layer writes to the slot of the cue $c$ that
ends at $p - 1$. The note $v_p$ and the write strength $\beta_p$ are computed from the hidden
state at $p$:
\begin{equation}
v_p = W_v x_p, \qquad \beta_p = \sigma(u^{\top} x_p + b), \qquad
S_c \leftarrow S_c + \beta_p\, (v_p - S_c) .
\label{eq:write}
\end{equation}
$W_v \in \mathbb{R}^{d_m \times d}$, $u \in \mathbb{R}^d$ and $b$ are learned for each cue length.
The note describes the token that followed the cue, and the slot is a weighted sum of the notes
written under the cue: its weights sum to less than one because the slot starts at zero, and of
two notes with equal write strength the later one has the larger weight.

At position $t$ and for each cue length, the layer reads the slot and the table row of the cue $c$
that ends at $t$:
\begin{equation}
r^N_t = R_c + S_c .
\label{eq:read}
\end{equation}
Here $S_c$ contains the writes of all positions before $t$. If the context has not yet written to
$c$, $S_c$ is zero and a flag $e^N_t$ is set to 1. A position $t < N$ has no cue of length $N$;
its read is zero with the flag set. Let $m_t$ be the concatenation of $r^N_t$ and
$e^N_t$ over the six cue lengths. The layer adds a gated projection of $m_t$ to the hidden state:
\begin{equation}
\alpha_t = \sigma\!\left( \frac{\langle \mathrm{norm}(x_t),\, \mathrm{norm}(W_K m_t) \rangle}{\sqrt{d}} \right),
\qquad x_t \leftarrow x_t + \alpha_t\, W_O m_t ,
\label{eq:gate}
\end{equation}
where $\mathrm{norm}$ is RMS normalization with a learned gain and $W_K, W_O \in \mathbb{R}^{d
\times 6(d_m + 1)}$. The gate $\alpha_t$ is a scalar that compares the read with the hidden state,
and the model can use it to attenuate a read that does not fit the current context.

When the cue that ends at $t$ also ended at an earlier position $t'$, its slot contains a note
computed from $x_{t'+1}$, the hidden state at the token that followed. The read therefore supplies
position $t$, which predicts $w_{t+1}$, with a representation of the earlier continuations of its
last $N$ tokens, for six values of $N$ at once. The table row is trained over the whole corpus,
and it is the entire read for a cue that the context has not written. The example of
Figure~\ref{fig:overview}b shows the two roles: in the first 16,384 tokens of that held-out
trajectory, the cue \texttt{tmp\_} occurs 18 times and is followed by \texttt{folder} each time,
whereas in the first 16,384 tokens of 2,000 sampled training trajectories it occurs 385 times,
followed by \texttt{path} in 230 of them (59.7\%) and by \texttt{folder} in none. With the writes disabled the
layer reduces to a static $N$-gram table, which we use as a control (Section~\ref{sec:replace}).

\subsection{Training and decoding}
\label{sec:method-train}

Because the addresses depend on token ids alone, the positions that write to each slot, and the
state that each read receives, are known before the forward pass. For every batch we compute this
schedule on the CPU: the writes are sorted by slot and position, and each read receives a pointer
to the last write that precedes it. The write of Equation~\ref{eq:write} is an affine map of the
slot, $S \mapsto (1 - \beta) S + \beta v$, and a composition of affine maps is affine. A parallel
prefix scan within each slot \citep{hillis1986} thus yields the state after every write, at a
depth logarithmic in the largest number of writes to one slot. The scan computes the same states
as sequential execution of the writes, and gradients reach the notes, the write strengths, the
gate and the table rows.

During decoding, each new token hashes one cue per cue length, reads one slot and one table row,
and writes one slot, at a cost that is independent of the length of the context. The state of a
session is one vector of width $d_m$ for every cue it has written, and a token adds at most one
slot per cue length; Section~\ref{sec:serving} reports the measured counts.

\subsection{Configuration}
\label{sec:method-config}

Unless we state otherwise, a model with \ingram{} uses one configuration, which we fixed before
the confirmatory runs that use it. Its cue lengths were selected on seeds 0 and 1
(Appendix~\ref{sec:cues}), after an earlier confirmatory comparison with cue lengths 1 to 5
(Appendix~\ref{app:seeds}).
It has the six cue lengths above, $d_m = 64$, $2^{16}$ table rows per cue
length, and one placement after the first block. The table has $6 \times 2^{16} \times 64 =
25.2$M parameters, and at model width 256 the projections of the layer have 0.3M. The dose grid of
Section~\ref{sec:types} also trains models with two placements, after blocks 1 and 5, and with
four, after blocks 1, 3, 5 and 7. Each placement has its own slots and projections, and all
placements of a model share one table. Appendix~\ref{app:other} reports a variant, the scored
read, which keeps the last $K$ notes of a slot separately, each with a key, and weights them and
the table row by attention from the reading position.
